%%%
%%% Radical "⼩"
%%%
\section*{Radical 42: ``⼩'' (⺌、⺍)}\addcontentsline{toc}{section}{Radical 42: ⼩,⺌、⺍}\addcontentsline{loh}{figure}{\#\#\#\# 42: ⼩}

%%%%%%%%%% 小 %%%%%%%%%%
\subsection*{小}\addcontentsline{loh}{figure}{小}

\begin{Entry}{小}{3}{⼩}[Kangxi 42]
  \begin{Phonetics}{小}{xiao3}[][HSK 1,2]
    \definition*{s.}{Sobrenome: Xiao}
    \definition{adj.}{menor; pequeno; insignificante; pouco; volume, área, quantidade, intensidade, etc. não são grandes | jovem | expressões humildes, referindo"-se a si mesmo ou a pessoas ou coisas relacionadas a si mesmo | por um tempo; por um curto período; por um curto período de tempo | o mais novo; o último na ordem de antiguidade; em último lugar na classificação}
    \definition{pref.}{usado antes do sobrenome, nome, posição na família, etc.}
    \definition{s.}{os jovens; pessoas mais jovens | concubina}
  \synonymref{略}{lve4}
  \synonymref{浅}{qian3}
  \synonymref{稍}{shao1}
  \synonymref{微}{wei1}
  \synonymref{细}{xi4}
  \synonymref{纤}{xian1}
  \synonymref{幼}{you4}
  \synonymref{暂}{zan4}
  \antonymref{大}{da4}
  \antonymref{巨}{ju4}
  \antonymref{老}{lao3}
  \end{Phonetics}
\end{Entry}

\begin{Entry}{小人}{3,2}{⼩,⼈}
  \begin{Phonetics}{小人}{xiao3ren2}[][HSK 7-9]
    \definition[个]{s.}{pessoa mesquinha; homem pequeno; homem ordinário; uma pessoa de caráter desprezível | ``seu humilde servo''; pessoa de baixa posição social; eu (usado por um homem de posição social inferior ao falar com seus superiores) | na antiguidade, referia"-se a pessoas de baixa condição social; era frequentemente usado como autorreferência | vilão; pessoa desprezível; personagem vil; pessoas imorais}
  \antonymref{君子}{jun1zi3}
  \end{Phonetics}
\end{Entry}

\begin{Entry}{小于}{3,3}{⼩,⼆}
  \begin{Phonetics}{小于}{xiao3yu2}[][HSK 6]
    \definition{prep.}{menor que; menos que; indica que um número ou quantidade é menor que outro}
  \end{Phonetics}
\end{Entry}

\begin{Entry}{小小}{3,3}{⼩,⼩}
  \begin{Phonetics}{小小}{xiao3xiao3}
    \definition{adj.}{muito pouco | muito pequeno}
  \end{Phonetics}
\end{Entry}

\begin{Entry}{小丑}{3,4}{⼩,⼀}
  \begin{Phonetics}{小丑}{xiao3chou3}[][HSK 7-9]
    \definition[个,名,位,帮]{s.}{palhaço; bufão; patife; na ópera tradicional chinesa, o palhaço é uma metáfora para alguém que se comporta de maneira indigna e é hábil em bajular os outros | um miserável desprezível; referindo"-se a uma pessoa insignificante}
  \end{Phonetics}
\end{Entry}

\begin{Entry}{小区}{3,4}{⼩,⼖}
  \begin{Phonetics}{小区}{xiao3qu1}[][HSK 7-9]
    \definition[个]{s.}{bairro; conjunto habitacional; comunidade residencial; áreas residenciais relativamente independentes na cidade, com instalações de serviços bem equipadas}
  \synonymref{社区}{she4qu1}
  \end{Phonetics}
\end{Entry}

\begin{Entry}{小心}{3,4}{⼩,⼼}
  \begin{Phonetics}{小心}{xiao3xin5}[][HSK 2]
    \definition{adj.}{cuidadoso; atento; com cautela}
    \definition{v.}{ter cuidado; ser cauteloso; estar atento; tomar cuidado; prestar atenção}
  \end{Phonetics}
\end{Entry}

\begin{Entry}{小心翼翼}{3,4,17,17}{⼩,⼼,⽻,⽻}
  \begin{Phonetics}{小心翼翼}{xiao3xin1-yi4yi4}[][HSK 7-9]
    \definition{expr.}{cautelosamente; cuidadosamente; com muita cautela; com grande cautela; originalmente significando solene e reverente, agora é usado para descrever ações extremamente cautelosas, sem espaço para negligência}
  \antonymref{粗心大意}{cu1xin1-da4yi4}
  \end{Phonetics}
\end{Entry}

\begin{Entry}{小气}{3,4}{⼩,⽓}
  \begin{Phonetics}{小气}{xiao3qi5}[][HSK 7-9]
    \definition{adj.}{mesquinho; apertado; avarento; parcimonioso; valorizar excessivamente os próprios bens | mesquinho; pequeno; de mente fechada}
  \antonymref{大量}{da4liang4}
  \antonymref{大气}{da4qi4}
  \antonymref{大方}{da4fang5}
  \antonymref{慷慨}{kang1kai3}
  \end{Phonetics}
\end{Entry}

\begin{Entry}{小气鬼}{3,4,9}{⼩,⽓,⿁}
  \begin{Phonetics}{小气鬼}{xiao3qi4gui3}
    \definition{adj.}{avarento; mesquinho; mão"-de"-vaca; miserável; pão"-duro}
  \end{Phonetics}
\end{Entry}

\begin{Entry}{小白菜}{3,5,11}{⼩,⽩,⾋}
  \begin{Phonetics}{小白菜}{xiao3bai2cai4}
    \definition[棵,些]{s.}{repolho chinês (\emph{Brassica chinensis}) | uma variedade de repolho chinês; \emph{pak choi} | \emph{bok choy}}
  \seealsoref{青菜}{qing1cai4}
  \seealsoref{小白菜儿}{xiao3bai2cai4r5}
  \end{Phonetics}
\end{Entry}

\begin{Entry}{小白菜儿}{3,5,11,2}{⼩,⽩,⾋,⼉}
  \begin{Phonetics}{小白菜儿}{xiao3bai2cai4r5}
    \definition{s.}{uma variedade de repolho chinês; \emph{pak choi}}
  \end{Phonetics}
\end{Entry}

\begin{Entry}{小众}{3,6}{⼩,⼈}
  \begin{Phonetics}{小众}{xiao3 zhong4}
    \definition{s.}{uma elite; uma minoria; um pequeno círculo; um grupo relativamente pequeno de pessoas; um pequeno grupo (no contexto de 大众)}
  \seealsoref{大众}{da4zhong4}
  \end{Phonetics}
\end{Entry}

\begin{Entry}{小伙子}{3,6,3}{⼩,⼈,⼦}
  \begin{Phonetics}{小伙子}{xiao3huo3zi5}[][HSK 4]
    \definition[位]{s.}{rapaz jovem; jovem colega}
  \end{Phonetics}
\end{Entry}

\begin{Entry}{小吃}{3,6}{⼩,⼝}
  \begin{Phonetics}{小吃}{xiao3chi1}[][HSK 4]
    \definition[家]{s.}{lanche; petiscos; comida com especialidades locais, não muito para uma porção | prato frio; prato feito; cortes de frios na culinária ocidental | pratos pequenos e baratos; pratos simples em restaurantes com porções pequenas e preços baixos}
  \end{Phonetics}
\end{Entry}

\begin{Entry}{小曲}{3,6}{⼩,⽈}
  \begin{Phonetics}{小曲}{xiao3qu1}[][HSK 7-9]
    \definition[首,支,段]{s.}{levedura para fabricação de vinho de arroz ou fermentação de arroz glutinoso | balada; melodia | melodia folclórica | música popular}
  \end{Phonetics}
\end{Entry}

\begin{Entry}{小声}{3,7}{⼩,⼠}
  \begin{Phonetics}{小声}{xiao3sheng1}[][HSK 2]
    \definition{v.}{falar em voz baixa; falar baixinho; sussurar}
  \end{Phonetics}
\end{Entry}

\begin{Entry}{小时}{3,7}{⼩,⽇}
  \begin{Phonetics}{小时}{xiao3shi2}[][HSK 1]
    \definition{clas.}{hora; unidade de medida legal do tempo, 1 hora equivale a 60 minutos, é 1/24 de um dia}
    \definition[个,半,些]{s.}{hora; refere"-se a um período de uma hora}
  \seealsoref{点钟}{dian3zhong1}
  \synonymref{点钟}{dian3zhong1}
  \synonymref{时刻}{shi2ke4}
  \synonymref{钟头}{zhong1tou2}
  \end{Phonetics}
\end{Entry}

\begin{Entry}{小时候}{3,7,10}{⼩,⽇,⼈}
  \begin{Phonetics}{小时候}{xiao3shi2hou5}[][HSK 2]
    \definition{s.}{na infância; quando alguém era jovem; refere"-se à infância}
  \end{Phonetics}
\end{Entry}

\begin{Entry}{小麦}{3,7}{⼩,⿆}
  \begin{Phonetics}{小麦}{xiao3mai4}[][HSK 6]
    \definition[粒,公斤,吨,棵]{s.}{trigo}
  \end{Phonetics}
\end{Entry}

\begin{Entry}{小卒}{3,8}{⼩,⼗}
  \begin{Phonetics}{小卒}{xiao3zu2}[][HSK 7-9]
    \definition{s.}{soldado de infantaria |Figurativo: ninguém | Xadrez: peão | particular | homem comum; um homem sem importância; pessoas comuns | um mero peão}
  \seealsoref{小卒子}{xiao3zu2 zi3}
  \antonymref{大王}{da4wang2}
  \antonymref{大王}{dai4wang5}
  \end{Phonetics}
\end{Entry}

\begin{Entry}{小卒子}{3,8,3}{⼩,⼗,⼦}
  \begin{Phonetics}{小卒子}{xiao3zu2 zi3}
    \definition{s.}{peão}
  \end{Phonetics}
\end{Entry}

\begin{Entry}{小姐}{3,8}{⼩,⼥}
  \begin{Phonetics}{小姐}{xiao3jie3}[][HSK 1]
    \definition[位,名,个,些]{s.}{jovem senhora; anteriormente, era assim que se referiam às filhas de famílias ricas. | senhorita; título honorífico para mulheres jovens | (gíria) prostituta}
  \seealsoref{姑娘}{gu1niang5}
  \synonymref{姑娘}{gu1niang5}
  \synonymref{女士}{nv3shi4}
  \end{Phonetics}
\end{Entry}

\begin{Entry}{小学}{3,8}{⼩,⼦}
  \begin{Phonetics}{小学}{xiao3xue2}[][HSK 1]
    \definition[所]{s.}{escola primária (ou fundamental); escolas que oferecem ensino fundamental básico | estudos filológicos; antigamente, referia"-se ao estudo da escrita, da fonética e da exegese}
  \antonymref{大学}{da4xue2}
  \end{Phonetics}
\end{Entry}

\begin{Entry}{小学生}{3,8,5}{⼩,⼦,⽣}
  \begin{Phonetics}{小学生}{xiao3xue2sheng5}[][HSK 1]
    \definition[个,位,名,些]{s.}{aluno; estudante; estudante do sexo masculino (男); estudante do sexo feminino (女) | um aluno mais novo (do que os outros da sua turma) | (dialeto) um menino pequeno}
  \seealsoref{男}{nan2}
  \seealsoref{女}{nv3}
  \end{Phonetics}
\end{Entry}

\begin{Entry}{小朋友}{3,8,4}{⼩,⽉,⼜}
  \begin{Phonetics}{小朋友}{xiao3peng2you3}[][HSK 1]
    \definition[个,名,位]{s.}{criança; crianças; refere"-se a crianças e adolescentes | (termo usado por um adulto para se dirigir a uma criança) amiguinho; menino (ou menina); termo carinhoso para se referir a crianças e adolescentes}
  \synonymref{少儿}{shao4'er2}
  \synonymref{小孩儿}{xiao3hai2r5}
  \antonymref{成人}{cheng2ren2}
  \antonymref{大人}{da4ren2}
  \end{Phonetics}
\end{Entry}

\begin{Entry}{小狗}{3,8}{⼩,⽝}
  \begin{Phonetics}{小狗}{xiao3gou3}
    \definition{s.}{filhote de cachorro; cães de pequeno porte, geralmente filhotes ou raças miniatura, são mantidos principalmente como animais de estimação}
  \end{Phonetics}
\end{Entry}

\begin{Entry}{小组}{3,8}{⼩,⽷}
  \begin{Phonetics}{小组}{xiao3zu3}[][HSK 2]
    \definition[个,名,位]{s.}{grupo; um pequeno grupo de pessoas}
  \end{Phonetics}
\end{Entry}

\begin{Entry}{小贩}{3,8}{⼩,⾙}
  \begin{Phonetics}{小贩}{xiao3fan4}[][HSK 7-9]
    \definition[个,名,位,帮]{s.}{vendedor ambulante; vendedor; mascate; comerciantes com muito pouco capital}
  \end{Phonetics}
\end{Entry}

\begin{Entry}{小品}{3,9}{⼩,⼝}
  \begin{Phonetics}{小品}{xiao3pin3}[][HSK 7-9]
    \definition[个,部,篇]{s.}{criação literária ou artística curta e simples; ensaio; esboço; esquete | esboço curto; ato curto; uma criação literária ou artística breve e simples; originalmente referindo"-se a versões abreviadas das escrituras budistas, agora se refere a ensaios curtos ou outras formas concisas de expressão}
  \end{Phonetics}
\end{Entry}

\begin{Entry}{小型}{3,9}{⼩,⼟}
  \begin{Phonetics}{小型}{xiao3xing2}[][HSK 4]
    \definition{adj.}{de tamanho pequeno; em pequena escala; miniatura; tipo pequeno; tamanho de bolso; tipo compacto}
    \definition{s.}{Mediterrâneo: escunas, pequenos veleiros de pesca ou turismo | pequeno \emph{rover} lunar (duas pessoas)}
  \end{Phonetics}
\end{Entry}

\begin{Entry}{小孩儿}{3,9,2}{⼩,⼦,⼉}
  \begin{Phonetics}{小孩儿}{xiao3hai2r5}[][HSK 1]
    \definition[个,名,位]{s.}{criança; bebê}
  \synonymref{儿童}{er2tong2}
  \synonymref{儿子}{er2zi5}
  \synonymref{孩子}{hai2zi5}
  \synonymref{女儿}{nv3'er2}
  \synonymref{小朋友}{xiao3peng2you3}
  \synonymref{子女}{zi3nv3}
  \antonymref{成人}{cheng2ren2}
  \antonymref{大人}{da4ren2}
  \end{Phonetics}
\end{Entry}

\begin{Entry}{小屋}{3,9}{⼩,⼫}
  \begin{Phonetics}{小屋}{xiao3wu1}
    \definition[间]{s.}{chalé; cabana; pousada | casita; cabana; casa geminada | cabine}
  \end{Phonetics}
\end{Entry}

\begin{Entry}{小树}{3,9}{⼩,⽊}
  \begin{Phonetics}{小树}{xiao3shu4}
    \definition[棵]{s.}{muda; árvore jovem; arbusto}
  \end{Phonetics}
\end{Entry}

\begin{Entry}{小洋白菜}{3,9,5,11}{⼩,⽔,⽩,⾋}
  \begin{Phonetics}{小洋白菜}{xiao3 yang2bai2cai4}
    \definition{s.}{couve de bruxelas}
  \end{Phonetics}
\end{Entry}

\begin{Entry}{小看}{3,9}{⼩,⽬}
  \begin{Phonetics}{小看}{xiao3kan4}[][HSK 7-9]
    \definition{v.}{subestimar; desprezar; menosprezar; desdenhar}
  \synonymref{鄙视}{bi3shi4}
  \synonymref{忽视}{hu1shi4}
  \synonymref{看不起}{kan4bu5qi3}
  \synonymref{歧视}{qi2shi4}
  \synonymref{无视}{wu2shi4}
  \antonymref{重视}{zhong4shi4}
  \end{Phonetics}
\end{Entry}

\begin{Entry}{小说}{3,9}{⼩,⾔}
  \begin{Phonetics}{小说}{xiao3shuo1}[][HSK 2]
    \definition[本,部,篇,章]{s.}{história; romance; ficção; uma forma literária que reflete a vida social por meio da descrição de personagens, ambiente e enredo}
  \end{Phonetics}
\end{Entry}

\begin{Entry}{小费}{3,9}{⼩,⾙}
  \begin{Phonetics}{小费}{xiao3fei4}[][HSK 6]
    \definition[笔]{s.}{gorjeta; gratificação; dinheiro extra pago por clientes e viajantes a funcionários de serviços em setores de serviços, como hotéis e pousadas}
  \end{Phonetics}
\end{Entry}

\begin{Entry}{小偷儿}{3,11,2}{⼩,⼈,⼉}
  \begin{Phonetics}{小偷儿}{xiao3tou1er5}[][HSK 5]
    \definition{s.}{ladrão insignificante (ou furtivo); ladrãozinho | ladrão}
  \end{Phonetics}
\end{Entry}

\begin{Entry}{小康}{3,11}{⼩,⼴}
  \begin{Phonetics}{小康}{xiao3kang1}[][HSK 7-9]
    \definition{adj.}{bem"-sucedido; (em relação ao padrão de vida) bem de vida; refere"-se à situação econômica de uma família que consegue manter um padrão de vida de classe média}
  \synonymref{富裕}{fu4yu4}
  \end{Phonetics}
\end{Entry}

\begin{Entry}{小提琴}{3,12,12}{⼩,⼿,⽟}
  \begin{Phonetics}{小提琴}{xiao3ti2qin2}[][HSK 7-9]
    \definition[个,些,把]{s.}{violino; um tipo de violino, o menor em tamanho e com o som mais agudo}
  \end{Phonetics}
\end{Entry}

\begin{Entry}{小溪}{3,13}{⼩,⽔}
  \begin{Phonetics}{小溪}{xiao3xi1}[][HSK 7-9]
    \definition[条]{s.}{riacho; córrego; um pequeno riacho raso, geralmente que flui por um vale ou floresta}
  \end{Phonetics}
\end{Entry}

\begin{Entry}{小腿}{3,13}{⼩,⾁}
  \begin{Phonetics}{小腿}{xiao3tui3}
    \definition{s.}{perna; parte inferior da perna (do joelho ao tornozelo)}
  \end{Phonetics}
\end{Entry}

\begin{Entry}{小路}{3,13}{⼩,⾜}
  \begin{Phonetics}{小路}{xiao3lu4}[][HSK 7-9]
    \definition[条]{s.}{caminho; trilha | atalho | faixa | estrada secundária | via}
  \antonymref{大道}{da4dao4}
  \end{Phonetics}
\end{Entry}

%%%%%%%%%% 少 %%%%%%%%%%
\subsection*{少}\addcontentsline{loh}{figure}{少}

\begin{Entry}{少}{4}{⼩}
  \begin{Phonetics}{少}{shao3}[][HSK 1]
    \definition{adj.}{menos; pouco; escasso; não atingir a quantidade original ou esperada}
    \definition{adv.}{um momento; um instante; provisoriamente; ligeiramente}
    \definition{v.}{faltar; ser insuficiente | dever | perder; desaparecer; extraviar | parar; desistir}
  \seealsoref{差}{cha1}
  \seealsoref{差}{cha4}
  \seealsoref{差}{chai1}
  \synonymref{差}{cha1}
  \synonymref{差}{cha4}
  \synonymref{差}{chai1}
  \synonymref{丢}{diu1}
  \synonymref{寡}{gua3}
  \synonymref{罕}{han3}
  \synonymref{缺}{que1}
  \synonymref{稍}{shao1}
  \synonymref{稀}{xi1}
  \synonymref{遗}{yi2}
  \antonymref{多}{duo1}
  \end{Phonetics}
  \begin{Phonetics}{少}{shao4}
    \definition*{s.}{Sobrenome: Shao}
    \definition{s.}{jovem}
    \definition{s.}{jovem mestre; filho de uma família rica}
  \antonymref{老}{lao3}
  \end{Phonetics}
\end{Entry}

\begin{Entry}{少儿}{4,2}{⼩,⼉}
  \begin{Phonetics}{少儿}{shao4'er2}[][HSK 6]
    \definition{s.}{criança}
  \end{Phonetics}
\end{Entry}

\begin{Entry}{少女}{4,3}{⼩,⼥}
  \begin{Phonetics}{少女}{shao4nv3}[][HSK 7-9]
    \definition[位,名]{s.}{empregada doméstica; jovem; moça; mulheres jovens solteiras}
  \end{Phonetics}
\end{Entry}

\begin{Entry}{少不了}{4,4,2}{⼩,⼀,⼅}
  \begin{Phonetics}{少不了}{shao3bu5liao3}[][HSK 7-9]
    \definition{v.}{não poder prescindir de; não poder dispensar | estar obrigado a; ser inevitável | não poder ser apenas alguns (ou pouco); não poder ser menos}
  \end{Phonetics}
\end{Entry}

\begin{Entry}{少见}{4,4}{⼩,⾒}
  \begin{Phonetics}{少见}{shao3jian4}[][HSK 7-9]
    \definition{adj.}{raramente visto; infrequente; raro}
    \definition{expr.}{``Não te vejo há muito tempo.'', ``Tenho te visto muito pouco ultimamente.'' ou ``Estou muito feliz em te ver novamente.''}
    \definition{v.}{difícil de ver | não ser familiar (para o falante) | ser raro (algo)}
  \end{Phonetics}
\end{Entry}

\begin{Entry}{少年}{4,6}{⼩,⼲}
  \begin{Phonetics}{少年}{shao4nian2}[][HSK 2]
    \definition[个,名,位]{s.}{adolescente; juventude; atualmente, a faixa etária geralmente referida é de 10 anos ou mais a 18 anos ou mais | menor; jovem; juvenil; refere"-se a menores na faixa etária anterior | jovem; adolescente; rapaz}
  \end{Phonetics}
\end{Entry}

\begin{Entry}{少有}{4,6}{⼩,⽉}
  \begin{Phonetics}{少有}{shao3you3}[][HSK 7-9]
    \definition{adj.}{raro; excepcional; escasso}
    \definition{adv.}{raramente}
  \end{Phonetics}
\end{Entry}

\begin{Entry}{少林寺}{4,8,6}{⼩,⽊,⼨}
  \begin{Phonetics}{少林寺}{shao4lin2si4}[][HSK 7-9]
    \definition*{s.}{Templo (ou Mosteiro) Shaolin, ao pé do Monte Song (嵩山) na província de Henan, onde o kung"-fu Shaolin foi desenvolvido}
  \seealsoref{嵩山}{song1shan1}
  \end{Phonetics}
\end{Entry}

\begin{Entry}{少量}{4,12}{⼩,⾥}
  \begin{Phonetics}{少量}{shao3liang4}[][HSK 7-9]
    \definition{s.}{uma pequena quantidade; um pouco; alguns; picada de pulga; toque; gosto; tapa; estalo; ninharia; gole}
  \end{Phonetics}
\end{Entry}

\begin{Entry}{少数}{4,13}{⼩,⽁}
  \begin{Phonetics}{少数}{shao3shu4}[][HSK 2]
    \definition{s.}{número pequeno; poucos; minoria}
  \end{Phonetics}
\end{Entry}

%%%%%%%%%% 尖 %%%%%%%%%%
\subsection*{尖}\addcontentsline{loh}{figure}{尖}

\begin{Entry}{尖}{6}{⼩}
  \begin{Phonetics}{尖}{jian1}[][HSK 6]
    \definition{adj.}{pontiagudo; afilado; agudo | agudo; estridente; penetrante | mesquinho; pão"-duro | mordaz; cáustico}
    \definition{s.}{ponto; ponta; topo | o melhor do seu tipo; a melhor escolha; a nata da safra; uma pessoa ou coisa notável}
    \definition{v.}{tornar (a voz, etc.) aguda; estridente}
  \end{Phonetics}
\end{Entry}

\begin{Entry}{尖锐}{6,12}{⼩,⾦}
  \begin{Phonetics}{尖锐}{jian1rui4}[][HSK 7-9]
    \definition{adj.}{pontiagudo; a ponta do objeto é fina e pequena, podendo facilmente entrar ou passar por outros objetos | incisivo; penetrante; conhecer e observar as coisas com muita precisão e profundidade | estridente; penetrante; o som é alto e fino, fazendo com que as pessoas se sintam desconfortáveis | afiado; intenso; muito intenso (discussão, luta, etc.)}
  \end{Phonetics}
\end{Entry}

\begin{Entry}{尖端}{6,14}{⼩,⽴}
  \begin{Phonetics}{尖端}{jian1duan1}[][HSK 7-9]
    \definition{adj.}{mais avançado; sofisticado; o mais alto nível de desenvolvimento (ciência e tecnologia, etc.)}
    \definition{s.}{ponta pontiaguda; pico; a ponta fina e pontiaguda de algo | topo; ápice; pico; uma metáfora para o mais alto nível de ciência e tecnologia}
  \end{Phonetics}
\end{Entry}

%%%%%%%%%% 尚 %%%%%%%%%%
\subsection*{尚}\addcontentsline{loh}{figure}{尚}

\begin{Entry}{尚}{8}{⼩}
  \begin{Phonetics}{尚}{shang4}[][HSK 7-9]
    \definition*{s.}{Sobrenome: Shang}
    \definition{adv.}{ainda}
    \definition{s.}{costume predominante; refere"-se à tendência predominante na sociedade; coisas que geralmente são admiradas pelas pessoas}
    \definition{v.}{valorizar; estimar; dar grande importância a; respeitar}
  \end{Phonetics}
\end{Entry}

\begin{Entry}{尚且}{8,5}{⼩,⼀}
  \begin{Phonetics}{尚且}{shang4qie3}
    \definition{conj.}{nem\dots; muito menos\dots; é usado antes do verbo da primeira oração de uma frase complexa para apresentar alguns exemplos óbvios para comparação, a segunda oração frequentemente usa 何况 ou 更 para ecoar e tirar conclusões inevitáveis sobre exemplos semelhantes com diferentes graus de gravidade}
  \seealsoref{更}{geng4}
  \seealsoref{何况}{he2kuang4}
  \end{Phonetics}
\end{Entry}

\begin{Entry}{尚且……何况……}{8,5,7,7}{⼩,⼀,⼈,⼎}
  \begin{Phonetics}{尚且……何况……}{shang4qie3 he2kuang4}
    \definition{conj.}{ainda que\dots, \dots; além do mais\dots e muito menos\dots}
  \end{Phonetics}
\end{Entry}

\begin{Entry}{尚未}{8,5}{⼩,⽊}
  \begin{Phonetics}{尚未}{shang4wei4}[][HSK 7-9]
    \definition{adv.}{ainda não}[问题尚未解决。===O problema ainda não foi resolvido.]
  \end{Phonetics}
\end{Entry}

%%%%%%%%%% 尝 %%%%%%%%%%
\subsection*{尝}\addcontentsline{loh}{figure}{尝}

\begin{Entry}{尝}{9}{⼩}
  \begin{Phonetics}{尝}{chang2}[][HSK 5]
    \definition{adv.}{alguma vez; uma vez}
    \definition{v.}{provar; experimentar o sabor de | provar; experimentar; conhecer | tentar; testar}
  \end{Phonetics}
\end{Entry}

\begin{Entry}{尝试}{9,8}{⼩,⾔}
  \begin{Phonetics}{尝试}{chang2shi4}[][HSK 5]
    \definition{v.}{tentar; provar; experimentar}
  \end{Phonetics}
\end{Entry}

%%%%% EOF %%%%%

